\documentclass[11pt,a4paper]{article}

\usepackage[a4paper,margin=1.65cm,top=1.45cm,bottom=1.55cm]{geometry}
\usepackage{fontspec}
\usepackage[french]{babel}
\usepackage{microtype}
\usepackage{array}
\usepackage{tabularx}
\usepackage{booktabs}
\usepackage{ragged2e}
\usepackage[table]{xcolor}
\usepackage{enumitem}
\usepackage{multicol}
\usepackage{titlesec}
\usepackage{fancyhdr}
\usepackage[most]{tcolorbox}
\usepackage{hyperref}
\usepackage{xurl}

\setmainfont{Arial}
\setsansfont{Arial}
\renewcommand{\familydefault}{\sfdefault}

\definecolor{primary}{HTML}{17365D}
\definecolor{secondary}{HTML}{2F75B5}
\definecolor{softblue}{HTML}{EAF3F8}
\definecolor{softgray}{HTML}{F3F5F7}
\definecolor{linegray}{HTML}{D5DCE3}
\definecolor{textgray}{HTML}{4C5967}
\definecolor{softgreen}{HTML}{EAF5EE}
\definecolor{softyellow}{HTML}{FFF7E6}
\definecolor{warning}{HTML}{A15C00}

\hypersetup{
  colorlinks=true,
  linkcolor=primary,
  urlcolor=secondary,
  pdftitle={Architecture offline-first et deux stacks simples - logiciel de caisse},
  pdfauthor={AMH Solution}
}
\urlstyle{same}

\setlength{\parindent}{0pt}
\setlength{\parskip}{4pt}
\renewcommand{\arraystretch}{1.24}
\setlist[itemize]{leftmargin=5mm,itemsep=1.5pt,topsep=2pt}
\setlist[enumerate]{leftmargin=6mm,itemsep=1.5pt,topsep=2pt}

\titleformat{\section}
  {\Large\bfseries\color{primary}}
  {\thesection.}{0.45em}{}
\titlespacing*{\section}{0pt}{10pt}{4pt}

\titleformat{\subsection}
  {\normalsize\bfseries\color{secondary}}
  {}{0pt}{}
\titlespacing*{\subsection}{0pt}{7pt}{3pt}

\pagestyle{fancy}
\setlength{\headheight}{22pt}
\fancyhf{}
\fancyhead[L]{\small\color{textgray}AMH Solution --- logiciel de caisse}
\fancyhead[R]{\small\color{textgray}Proposition technique V1}
\fancyfoot[C]{\small\color{textgray}\thepage}
\renewcommand{\headrulewidth}{0.35pt}
\renewcommand{\headrule}{\hbox to\headwidth{\color{linegray}\leaders\hrule height \headrulewidth\hfill}}

\newcolumntype{Y}{>{\RaggedRight\arraybackslash}X}
\newcolumntype{L}[1]{>{\RaggedRight\arraybackslash}p{#1}}
\newcolumntype{C}[1]{>{\Centering\arraybackslash}p{#1}}

\newcommand{\yes}{\textbf{\color{primary}Oui}}
\newcommand{\later}{\textbf{\color{textgray}Plus tard}}
\newcommand{\source}[2]{\href{#2}{#1}}

\begin{document}

\begin{center}
  {\fontsize{22}{26}\selectfont\bfseries\color{primary}
  Architecture offline-first et deux stacks simples}\\[3pt]
  {\large Logiciel de caisse pour restaurant}\\[7pt]
  {\small\color{textgray}Note historique --- version 1 --- 2 juin 2026}
\end{center}

\begin{tcolorbox}[
  colback=softyellow,
  colframe=warning,
  boxrule=0.8pt,
  arc=2pt,
  left=8pt,right=8pt,top=7pt,bottom=7pt
]
\textbf{\color{warning}Document remplacé pour la V1}\\
Cette note comparait initialement des options desktop Windows. La décision du 5 juin 2026 retient
désormais une application native \textbf{Android Kotlin / Jetpack Compose} livrée sous forme
d'\textbf{APK} pour tablette Android ou terminal POS Android. Ce document reste conservé comme
trace de réflexion, mais il n'est plus la référence de développement de la V1.
\end{tcolorbox}

\begin{tcolorbox}[
  colback=softgreen,
  colframe=primary,
  boxrule=0.8pt,
  arc=2pt,
  left=8pt,right=8pt,top=7pt,bottom=7pt
]
\textbf{\color{primary}Recommandation principale}\\
Retenir un \textbf{client épais local} sur un \textbf{POS ou une tablette Windows} et comparer
deux solutions.
Pour livrer une première version exploitable avec le moins de complexité, privilégier
\textbf{Python + PySide6 + SQLite}. Conserver \textbf{Electron + React + TypeScript + SQLite}
comme alternative si l'apprentissage d'une application desktop basée sur les technologies web
est prioritaire. Dans les deux cas, une vente reste possible sans Internet.
\end{tcolorbox}

\section{Objectif et principe}

Le logiciel doit rester simple à développer, installer et maintenir. La première version suit un
principe \textbf{offline-first} : elle fonctionne sans connexion Internet sur un seul appareil
Windows. L'objectif du stage est aussi de progresser en logique métier, en SQL et, selon la stack
retenue, en Python ou en TypeScript.

\begin{itemize}
  \item Une vente est enregistrée dans SQLite sur la caisse avant toute synchronisation éventuelle.
  \item Le logiciel doit continuer à fonctionner si Internet est indisponible.
  \item Une sauvegarde manuelle ou planifiée doit pouvoir être exportée sur un autre support.
  \item Le choix du matériel sera fixé après un test réel d'impression.
\end{itemize}

\begin{tcolorbox}[
  colback=softblue,
  colframe=secondary,
  boxrule=0.55pt,
  arc=2pt,
  left=8pt,right=8pt,top=6pt,bottom=6pt
]
\textbf{Pourquoi cibler Windows dans la première version ?}\\
Python s'exécute et se transforme simplement en application installable sur Windows. Sur Android
ou iPadOS, intégrer Python dans une application locale ajoute des difficultés qui n'apportent pas
de valeur au premier prototype. Mieux vaut réussir un produit simple sur un OS avant d'envisager
un portage.
\end{tcolorbox}

\section{Deux stacks techniques minimalistes}

Les deux propositions ciblent Windows et partagent la même base : \textbf{SQLite} comme source de
vérité, \textbf{JSON} pour les tickets figés, \textbf{CSV} pour les exports et \textbf{Git} pour le
suivi du code. Les montants sont stockés en centimes sous forme d'entiers, par exemple
\texttt{1850} pour 18,50 euros.

\subsection{Option A --- application desktop Python : choix recommandé}

\rowcolors{2}{softgray}{white}
\begin{tabularx}{\textwidth}{L{3.15cm} L{3.15cm} Y}
  \toprule
  \rowcolor{primary}
  \color{white}\textbf{Besoin} &
  \color{white}\textbf{Choix} &
  \color{white}\textbf{Rôle} \\
  \midrule
  Logique métier &
  \textbf{Python} &
  Écrire les règles essentielles : panier, paiements, historique et statistiques simples. \\
  Interface tactile &
  \textbf{PySide6} &
  Construire une application desktop tactile avec Qt for Python. \\
  Stockage local &
  \textbf{SQLite} &
  Enregistrer les données dans un fichier local, sans serveur à installer. \\
  Tickets et sauvegardes &
  \textbf{JSON, CSV} &
  Figer les tickets et produire des exports lisibles. \\
  Exécutable Windows &
  \textbf{pyside6-deploy} &
  Générer un exécutable Windows à partir de l'application Python. \\
  Tests &
  \textbf{pytest} &
  Vérifier les calculs critiques et les annulations. \\
  Impression &
  SDK du fabricant &
  Piloter l'imprimante après un test sur le matériel réel. \\
  \bottomrule
\end{tabularx}
\rowcolors{2}{}{}

\vspace{3pt}
\textbf{Atout principal :} une seule base de code Python, sans API locale et sans build frontend.
C'est le chemin le plus court pour obtenir une application installable et se concentrer sur la
logique métier.

\subsection{Option B --- application desktop Electron et React}

\rowcolors{2}{softgray}{white}
\begin{tabularx}{\textwidth}{L{3.15cm} L{3.15cm} Y}
  \toprule
  \rowcolor{primary}
  \color{white}\textbf{Besoin} &
  \color{white}\textbf{Choix} &
  \color{white}\textbf{Rôle} \\
  \midrule
  Runtime desktop &
  \textbf{Electron} &
  Exécuter une application desktop locale qui embarque Chromium et Node.js. \\
  Logique métier &
  \textbf{TypeScript} &
  Écrire les règles essentielles dans un langage typé. \\
  Interface tactile &
  \textbf{React, TypeScript, CSS} &
  Construire des écrans tactiles structurés en composants réutilisables. \\
  Build interface &
  \textbf{Vite} &
  Transformer l'interface React en fichiers statiques intégrés à l'application. \\
  Données et matériel &
  Processus principal Electron via \textbf{IPC} &
  Relier l'interface à SQLite et aux fonctions natives du poste. \\
  Stockage local &
  \textbf{SQLite} &
  Enregistrer les données dans un fichier local, sans serveur distant. \\
  Tickets et sauvegardes &
  \textbf{JSON, CSV} &
  Figer les tickets et produire des exports lisibles. \\
  Exécutable Windows &
  \textbf{Electron Forge} &
  Générer un exécutable ou un installateur Windows. \\
  Tests &
  \textbf{Vitest} &
  Vérifier les calculs critiques et les annulations. \\
  \bottomrule
\end{tabularx}
\rowcolors{2}{}{}

\vspace{3pt}
\textbf{Atout principal :} progresser en React et TypeScript dans une vraie application desktop.
Electron permet de conserver une interface web moderne tout en accédant aux fonctions locales
par IPC. Cette option reste raisonnable, mais comporte davantage de pièces à assembler.

\subsection{Comparaison rapide}

\rowcolors{2}{softgray}{white}
\begin{tabularx}{\textwidth}{L{4.25cm} Y Y}
  \toprule
  \rowcolor{primary}
  \color{white}\textbf{Critère} &
  \color{white}\textbf{Option A --- PySide6} &
  \color{white}\textbf{Option B --- Electron + React} \\
  \midrule
  Simplicité de développement &
  \textbf{Très bonne} &
  Bonne \\
  Nombre de briques principales &
  Python, PySide6, SQLite &
  Electron, React, TypeScript, Vite, SQLite \\
  Compétences travaillées &
  Python, SQL, logique métier &
  TypeScript, SQL, React, IPC, application desktop \\
  Priorité &
  \textbf{Livrer efficacement} &
  Approfondir le desktop avec technologies web \\
  \bottomrule
\end{tabularx}
\rowcolors{2}{}{}

\section{Principe offline-first et évolution possible}

Le site public de Splash présente une caisse connectée, un mode hors ligne fiable, des statistiques,
un back-office, des bornes et un écran cuisine. Le détail de leur code n'est pas public. Le schéma
ci-dessous est donc une \textbf{trajectoire prudente pour ce projet}, pas une description de leur
implémentation interne.

\rowcolors{2}{softgray}{white}
\begin{tabularx}{\textwidth}{L{4.1cm} Y Y}
  \toprule
  \rowcolor{primary}
  \color{white}\textbf{Sujet} &
  \color{white}\textbf{Première version à livrer} &
  \color{white}\textbf{Extension facultative ultérieure} \\
  \midrule
  Encaissement &
  Application desktop et SQLite locale &
  Identique : le cloud ne bloque jamais une vente \\
  Données &
  Catalogue, commandes et paiements locaux &
  File de synchronisation envoyée quand Internet revient \\
  Pilotage &
  Historique et statistiques sur la caisse &
  API centrale, base PostgreSQL et back-office web \\
  Périphériques &
  Imprimante et TPE testés sur le poste &
  Borne ou écran cuisine ajoutés selon le besoin \\
  \bottomrule
\end{tabularx}
\rowcolors{2}{}{}

\vspace{4pt}
\begin{tcolorbox}[
  colback=softblue,
  colframe=secondary,
  boxrule=0.55pt,
  arc=2pt,
  left=8pt,right=8pt,top=6pt,bottom=6pt
]
\textbf{Place éventuelle de Next.js}\\
Next.js peut devenir utile plus tard pour un back-office web consultable à distance. Il ne doit
pas être l'unique runtime de la caisse au moment de l'encaissement.
\end{tcolorbox}

\section{Compatibilité selon le système}

Les deux stacks proposées ciblent Windows. Les autres systèmes restent intéressants à comparer
pour le matériel, mais demanderaient un autre choix technique si un portage devenait nécessaire.

\rowcolors{2}{softgray}{white}
\begin{tabularx}{\textwidth}{L{2.45cm} L{3.15cm} L{3.35cm} Y}
  \toprule
  \rowcolor{primary}
  \color{white}\textbf{Système} &
  \color{white}\textbf{Choix technique} &
  \color{white}\textbf{Langages principaux} &
  \color{white}\textbf{Remarque} \\
  \midrule
  Android / SUNMI OS &
  Non retenu pour la V1 &
  Kotlin ou JavaScript &
  Bon matériel POS, mais Python local n'est pas le chemin le plus simple. \\
  Windows 11 &
  \textbf{Choix recommandé} &
  Python ou TypeScript &
  Simple à développer, empaqueter et tester ; compatible avec de nombreux périphériques USB. \\
  iPadOS &
  Non retenu pour la V1 &
  Swift ou JavaScript &
  Bon matériel, mais distribution et impression plus contraignantes pour cette première version. \\
  \bottomrule
\end{tabularx}
\rowcolors{2}{}{}

\section{Choix conseillé pour commencer}

\begin{tcolorbox}[
  colback=softblue,
  colframe=secondary,
  boxrule=0.55pt,
  arc=2pt,
  left=8pt,right=8pt,top=7pt,bottom=7pt
]
\textbf{\color{primary}Premier environnement à tester : un appareil Windows 11 disponible}\\
Avant tout achat, construire le prototype desktop sur un ordinateur ou une tablette Windows déjà
disponible. Tester ensuite le logiciel sur un terminal POS Windows, par exemple un SUNMI T3 PRO
SUPER configuré sous Windows si cette version est disponible auprès du revendeur.
\end{tcolorbox}

\textbf{Pourquoi procéder ainsi :} le choix définitif du terminal dépend surtout de l'impression,
du tiroir-caisse, du prix et de la disponibilité. Le prototype peut avancer sans attendre cet achat.

\section{Plan d'action détaillé}

L'objectif est de valider chaque bloc avant de passer au suivant. Ce plan donne un ordre de travail,
pas un calendrier figé.

\begin{enumerate}
  \item \textbf{Valider les décisions de départ.}\\
  Choisir l'option A ou B avec le maître de stage, préciser les moyens de paiement, les taux de TVA,
  le besoin d'impression et le matériel déjà disponible. \textbf{Résultat :} une décision technique
  courte et les questions client prioritaires.

  \item \textbf{Créer le socle de l'application desktop.}\\
  Initialiser le dépôt Git, l'arborescence, la configuration et un premier écran tactile lançable.
  \textbf{Résultat :} une application vide mais installable en développement.

  \item \textbf{Modéliser les données locales.}\\
  Créer SQLite avec articles, catégories, commandes, lignes de commande, paiements, utilisateurs
  et identifiants stables. Stocker les montants en centimes. \textbf{Résultat :} un schéma versionné
  et quelques données de démonstration.

  \item \textbf{Développer catalogue, panier et commandes.}\\
  Ajouter les articles au panier, modifier les quantités et distinguer sur place, à emporter et
  livraison. \textbf{Résultat :} une commande complète peut être préparée avant encaissement.

  \item \textbf{Implémenter les encaissements.}\\
  Gérer espèces, carte, titres-restaurant, autres moyens, paiement mixte et rendu de monnaie.
  \textbf{Résultat :} les calculs critiques sont couverts par des tests automatisés.

  \item \textbf{Ajouter la traçabilité.}\\
  Figer un instantané JSON à l'encaissement, conserver l'opérateur et l'horodatage, puis traiter
  une correction ou une annulation comme une nouvelle opération liée. \textbf{Résultat :} aucune
  vente validée n'est modifiée silencieusement.

  \item \textbf{Construire l'historique et les statistiques simples.}\\
  Rechercher les ventes et afficher chiffre d'affaires, volume de transactions, répartition des
  paiements et totaux journaliers. \textbf{Résultat :} le responsable dispose d'un premier tableau
  de bord local.

  \item \textbf{Tester l'impression et les périphériques.}\\
  Imprimer une note sur le matériel réel, gérer l'échec d'impression et valider le besoin de
  tiroir-caisse. \textbf{Résultat :} un compte rendu matériel avant achat définitif.

  \item \textbf{Ajouter sauvegarde, restauration et livraison Windows.}\\
  Exporter JSON et CSV, tester la restauration, puis générer l'exécutable ou l'installateur.
  \textbf{Résultat :} une version démontrable hors ligne sur un poste Windows distinct.

  \item \textbf{Réaliser la recette et la mise en service.}\\
  Jouer les parcours essentiels avec le client, corriger les défauts bloquants, rédiger un guide
  court et former les utilisateurs. \textbf{Résultat :} une première version exploitable validée.
\end{enumerate}

\section{Décisions à prendre avant achat}

\begin{itemize}
  \item La caisse reste-t-elle fixée au comptoir ou doit-elle être déplacée en salle ?
  \item Une imprimante ticket 80 mm intégrée est-elle souhaitée ?
  \item Faut-il ouvrir automatiquement un tiroir-caisse ?
  \item Une imprimante cuisine est-elle nécessaire dès la première version ?
  \item Quel budget matériel est prévu ?
  \item Le client possède-t-il déjà une tablette ou un terminal réutilisable ?
\end{itemize}

\clearpage
\section{Comparaison du matériel et ordre de prix}

Les modèles ci-dessous sont des pistes de comparaison, pas une commande d'achat. Les prix sont des
\textbf{ordres de grandeur TTC relevés le 2 juin 2026}, hors accessoires et hors promotions
temporaires lorsqu'elles sont identifiables. Ils doivent être revérifiés avant achat. Pour un
terminal professionnel sans tarif public stable en France, il faut demander un devis au revendeur.

\subsection{Android et terminaux POS Android}

\rowcolors{2}{softgray}{white}
\begin{tabularx}{\textwidth}{L{3.2cm} L{2.45cm} Y L{3.35cm}}
  \toprule
  \rowcolor{primary}
  \color{white}\textbf{Modèle} &
  \color{white}\textbf{Prix indicatif} &
  \color{white}\textbf{Intérêt} &
  \color{white}\textbf{Vigilance} \\
  \midrule
  Samsung Galaxy Tab A9+ &
  Environ 250 à 350 EUR &
  Tablette 11 pouces économique avec imprimante externe. &
  Prévoir support, alimentation continue et imprimante réseau. \\
  Samsung Galaxy Tab Active5 Pro Enterprise Edition &
  Environ 600 à 850 EUR &
  Tablette renforcée 10,1 pouces ; protection IP68 et batterie amovible. &
  Plus coûteuse ; imprimante externe nécessaire. \\
  SUNMI D3 MINI &
  Environ 400 à 500 EUR &
  POS Android compact 10,1 pouces avec imprimante thermique 80 mm intégrée. &
  Ne correspond pas aux stacks Windows retenues pour la V1. \\
  \bottomrule
\end{tabularx}
\rowcolors{2}{}{}

\subsection{Windows}

\rowcolors{2}{softgray}{white}
\begin{tabularx}{\textwidth}{L{3.2cm} L{2.45cm} Y L{3.35cm}}
  \toprule
  \rowcolor{primary}
  \color{white}\textbf{Modèle} &
  \color{white}\textbf{Prix indicatif} &
  \color{white}\textbf{Intérêt} &
  \color{white}\textbf{Vigilance} \\
  \midrule
  Microsoft Surface Pro 12 ou 13 pouces &
  Environ 1 000 à 1 600 EUR &
  Tablette Windows polyvalente pour le développement et les essais. &
  Ajouter support et imprimante ; vérifier les pilotes périphériques sur processeur ARM. \\
  Lenovo ThinkPad X12 Detachable Gen 2 &
  Environ 1 600 à 2 700 EUR &
  Tablette Windows professionnelle 12 pouces avec clavier détachable. &
  Robuste mais chère pour une caisse seule ; comparer avec un terminal POS dédié. \\
  SUNMI T3 PRO SUPER &
  Sur devis &
  POS fixe dual-OS Windows ou Android, adapté à un comptoir. &
  Valider la configuration Windows disponible, l'imprimante et les accessoires. \\
  \bottomrule
\end{tabularx}
\rowcolors{2}{}{}

\subsection{iPadOS}

\rowcolors{2}{softgray}{white}
\begin{tabularx}{\textwidth}{L{3.2cm} L{2.45cm} Y L{3.35cm}}
  \toprule
  \rowcolor{primary}
  \color{white}\textbf{Modèle} &
  \color{white}\textbf{Prix indicatif} &
  \color{white}\textbf{Intérêt} &
  \color{white}\textbf{Vigilance} \\
  \midrule
  iPad 11 pouces (A16) &
  Environ 400 à 650 EUR &
  Option iPad la plus raisonnable pour une caisse simple. &
  Prévoir support, alimentation et validation de l'impression. \\
  iPad Air 11 pouces (M4) &
  Environ 650 à 1 100 EUR &
  Plus confortable si le client veut un appareil polyvalent. &
  Plus puissant et plus cher que nécessaire pour la caisse seule. \\
  iPad mini (A17 Pro) &
  Environ 550 à 900 EUR &
  Intéressant pour la mobilité ou la prise de commande en salle. &
  Écran 8,3 pouces trop petit pour un comptoir fixe. \\
  \bottomrule
\end{tabularx}
\rowcolors{2}{}{}

\clearpage
\section{Sources consultées}

{\footnotesize
\begin{multicols}{2}
\begin{itemize}[itemsep=0pt,topsep=0pt,leftmargin=4mm]
  \item Python \texttt{sqlite3} :
  \source{documentation}{https://docs.python.org/3/library/sqlite3.html}
  \item PySide6 :
  \source{documentation}{https://doc.qt.io/qtforpython-6/gettingstarted.html}
  \item \texttt{pyside6-deploy} :
  \source{documentation}{https://doc.qt.io/qtforpython-6/deployment/deployment-pyside6-deploy.html}
  \item Splash :
  \source{site officiel}{https://splash360.fr/}
  \item Electron :
  \source{documentation}{https://www.electronjs.org/docs/latest/tutorial/tutorial-prerequisites}
  \item Electron : processus principal et rendu :
  \source{documentation}{https://www.electronjs.org/docs/latest/tutorial/process-model}
  \item Electron Forge :
  \source{documentation}{https://www.electronforge.io/}
  \item React :
  \source{documentation}{https://react.dev/learn/build-a-react-app-from-scratch}
  \item Vite :
  \source{documentation}{https://vite.dev/guide/}
  \item Vitest :
  \source{documentation}{https://vitest.dev/guide/}
  \item Next.js :
  \source{documentation}{https://nextjs.org/docs}
  \item SQLite :
  \source{présentation officielle}{https://sqlite.org/about.html}
  \item Samsung Galaxy Tab A9+ :
  \source{Samsung France}{https://www.samsung.com/fr/tablets/galaxy-tab-a/galaxy-tab-a9-plus-wifi-graphite-128gb-sm-x210nzaeeub/}
  \item Prix Galaxy Tab A9+ :
  \source{Fnac}{https://www.fnac.com/Tablette-Tactile-Samsung-Galaxy-Tab-A9-11-Wifi-128-Go-Anthracite/a21688745/w-4}
  \item Galaxy Tab Active5 Pro :
  \source{Samsung Professionnels France}{https://www.samsung.com/fr/business/tablets/galaxy-tab-active/galaxy-tab-active5-pro-enterprise-edition-sm-x356bzgaeeb/}
  \item Prix Galaxy Tab Active5 Pro :
  \source{Galaxus France}{https://www.galaxus.fr/fr/s1/product/samsung-galaxy-tab-active5-pro-edition-entreprise-5g-1010-128-go-vert-tablette-57860971}
  \item SUNMI D3 MINI :
  \source{fiche produit PDF}{https://file.cdn.sunmi.com/newebsite/products/d3-mini/specs/d3-mini-en.pdf}
  \item Prix SUNMI D3 MINI :
  \source{Shop NFC}{https://www.shopnfc.com/fr/terminaux-nfc-sunmi/679-sunmi-d3-mini-terminal-de-point-de-vente-android-gms-imprimante-80-mm-cable-ue.html}
  \item SUNMI T3 PRO :
  \source{SUNMI}{https://www.sunmi.com/en-US/t3-pro-series/}
  \item Microsoft Surface Pro :
  \source{Microsoft France}{https://www.microsoft.com/fr-fr/surface/devices/surface-pro/}
  \item Prix Surface Pro :
  \source{Microsoft Store France}{https://www.microsoft.com/fr-fr/store/collections/consumersurface/}
  \item Lenovo ThinkPad X12 Detachable Gen 2 :
  \source{Lenovo France}{https://www.lenovo.com/fr/fr/p/laptops/thinkpad/thinkpadx/lenovo-thinkpad-x12-gen-2-12-inch-intel/21lkcto1wwfr2}
  \item Prix ThinkPad X12 Gen 2 :
  \source{PcComponentes France}{https://www.pccomponentes.fr/portable-lenovo-thinkpad-x12-detachable-gen-2-123-intel-core-ultra-5-134u-16gb-512gb-ssd-intel-graphics-windows-11-pro}
  \item Apple iPad 11 pouces (A16) :
  \source{Apple}{https://www.apple.com/fr/ipad-11/}
  \item Apple iPad Air :
  \source{Apple}{https://www.apple.com/fr/ipad-air/}
  \item Apple iPad mini (A17 Pro) :
  \source{Apple}{https://support.apple.com/fr-fr/121456}
  \item Prix des iPad :
  \source{Apple Store France}{https://www.apple.com/fr/shop/buy-ipad}
\end{itemize}
\end{multicols}
}

\vfill
{\small\color{textgray}
Document de travail. Les modèles cités sont des candidats à tester. La décision d'achat dépend
de la disponibilité en France, du prix obtenu auprès du revendeur et de la validation de
l'impression sur le matériel réel.
}

\end{document}
